\section{长短期记忆网络（LSTM）}

\subsection{设计动机}

简单 RNN 的核心问题在于：隐藏状态在每个时间步被\textbf{完全改写}（通过矩阵乘法），这使得信息很难被保留。我们需要一种架构，使得信息可以被\textbf{选择性地添加和遗忘}，而非每步全部覆盖。

\begin{figure}[H]
\centering
\includegraphics[width=0.95\textwidth]{slides-images/slide-019.jpg}
\caption{RNN 信息保存的困难：隐藏状态在每步被矩阵乘法完全重写，无法有效保留信息\protect\footnotemark}
\end{figure}
\footnotetext{Slides 第20页。}

\begin{knowledgebox}{LSTM 的历史}
LSTM（Long Short-Term Memory）由 Hochreiter 和 Schmidhuber 于 1997 年提出。名称的含义是``长期的短期记忆''——人类的短期记忆可以保持相当长时间（几分钟到几十分钟），而简单 RNN 的``短期记忆''仅约 7 个 token。LSTM 旨在延长这一有效记忆窗口。

值得注意的是，现代常用 LSTM 的一个关键组件——遗忘门——并不在 1997 年的原始论文中，而是由 Gers 和 Schmidhuber 在 2000 年的后续工作中引入的。Schmidhuber 的学生 Alex Graves 后来将 LSTM 带到 Hinton 的实验室，最终经由 Google 在 2014--2016 年间大规模部署，LSTM 才成为主流。
\end{knowledgebox}

\subsection{LSTM 的核心组件}

LSTM 在每个时间步维护两个状态向量：

\begin{itemize}
    \item \textbf{隐藏状态}（Hidden State）$h^{(t)}$：用于预测输出和传递给下一步
    \item \textbf{细胞状态}（Cell State）$c^{(t)}$：作为长期记忆的载体
\end{itemize}

信息的流动由三个\textbf{门}（Gate）控制，每个门是一个与隐藏状态同维度的向量，值在 $[0, 1]$ 之间：

\begin{figure}[H]
\centering
\includegraphics[width=0.95\textwidth]{slides-images/slide-020.jpg}
\caption{LSTM 的核心概念：隐藏状态 $h^{(t)}$ 和细胞状态 $c^{(t)}$，以及三个门控机制\protect\footnotemark}
\end{figure}
\footnotetext{Slides 第21页。}

\subsection{LSTM 的数学公式}

\begin{figure}[H]
\centering
\includegraphics[width=0.95\textwidth]{slides-images/slide-021.jpg}
\caption{LSTM 的完整数学公式：三个门的计算、细胞状态更新和隐藏状态计算\protect\footnotemark}
\end{figure}
\footnotetext{Slides 第22页。}

\textbf{第一步：计算三个门}

$$f^{(t)} = \sigma\left(W_f h^{(t-1)} + U_f x^{(t)} + b_f\right)$$
$$i^{(t)} = \sigma\left(W_i h^{(t-1)} + U_i x^{(t)} + b_i\right)$$
$$o^{(t)} = \sigma\left(W_o h^{(t-1)} + U_o x^{(t)} + b_o\right)$$

各符号说明：
\begin{itemize}
    \item $f^{(t)}$：\textbf{遗忘门}（Forget Gate），控制保留多少旧的细胞状态。$\sigma$ 为 sigmoid 函数，输出在 $[0,1]$ 之间
    \item $i^{(t)}$：\textbf{输入门}（Input Gate），控制写入多少新信息到细胞
    \item $o^{(t)}$：\textbf{输出门}（Output Gate），控制从细胞中读出多少信息到隐藏状态
    \item $W_f, W_i, W_o$：隐藏状态到各门的权重矩阵
    \item $U_f, U_i, U_o$：输入到各门的权重矩阵
    \item $b_f, b_i, b_o$：偏置项
\end{itemize}

\textbf{第二步：计算候选细胞内容}

$$\tilde{c}^{(t)} = \tanh\left(W_c h^{(t-1)} + U_c x^{(t)} + b_c\right)$$

这个候选内容的计算形式与简单 RNN 完全相同，只是使用了专用的权重矩阵 $W_c$ 和 $U_c$。

\textbf{第三步：更新细胞状态}

$$c^{(t)} = f^{(t)} \odot c^{(t-1)} + i^{(t)} \odot \tilde{c}^{(t)}$$

\begin{itemize}
    \item $\odot$：Hadamard 积（逐元素乘法）
    \item $f^{(t)} \odot c^{(t-1)}$：遗忘门控制保留多少旧记忆
    \item $i^{(t)} \odot \tilde{c}^{(t)}$：输入门控制写入多少新信息
\end{itemize}

\textbf{第四步：计算隐藏状态}

$$h^{(t)} = o^{(t)} \odot \tanh(c^{(t)})$$

输出门控制从细胞状态中``读出''多少信息到隐藏状态。$\tanh$ 将细胞状态的值压缩到 $[-1, 1]$。

\begin{warningbox}{遗忘门实际上是``记忆门''}
Manning 指出，``遗忘门''的命名容易引起误解。当 $f^{(t)} = 1$ 时，完全\textbf{保留}旧记忆；当 $f^{(t)} = 0$ 时，完全\textbf{遗忘}。因此，将其理解为``记忆门''（Remember Gate）更符合直觉——它的值越大，保留得越多。
\end{warningbox}

\subsection{LSTM 的图示理解}

\begin{figure}[H]
\centering
\includegraphics[width=0.95\textwidth]{slides-images/slide-022.jpg}
\caption{LSTM 的 Olah 风格图示：展示信息在 LSTM 单元中的完整流动过程。图来自 Chris Olah 的博客\protect\footnotemark}
\end{figure}
\footnotetext{Slides 第23页。来源：\url{http://colah.github.io/posts/2015-08-Understanding-LSTMs/}}

\begin{figure}[H]
\centering
\includegraphics[width=0.95\textwidth]{slides-images/slide-023.jpg}
\caption{LSTM 内部结构带标注版：清晰标注了遗忘、输入、输出各步骤的数据流。注意右上角的``加号是秘密''提示\protect\footnotemark}
\end{figure}
\footnotetext{Slides 第24页。来源：Chris Olah 博客。}

\subsection{LSTM 为何有效：加法的秘密}

\begin{importantbox}{加法操作是 LSTM 的核心秘密}
简单 RNN 的隐藏状态更新是纯\textbf{乘法}的（$h^{(t)} = \sigma(W_h h^{(t-1)} + \ldots)$），信息在乘法操作中容易消失或爆炸。而 LSTM 的细胞状态更新包含一个\textbf{加法}操作：
$$c^{(t)} = f^{(t)} \odot c^{(t-1)} \textcolor{red}{+} \, i^{(t)} \odot \tilde{c}^{(t)}$$
这个加号是关键——它使得信息可以沿细胞状态直接流动，无需经过矩阵乘法。当 $f^{(t)} = 1$ 且 $i^{(t)} = 0$ 时，$c^{(t)} = c^{(t-1)}$，信息可以无损地线性传递。这从根本上解决了梯度消失问题。
\end{importantbox}

这种加法更新机制与人类记忆的工作方式类似：我们通过\textbf{添加}新记忆来更新认知，而非每次完全重写。最初的 LSTM（1997）甚至没有遗忘门，是纯粹的加法机制。但纯加法在长序列上会累积过多信息导致功能失调，因此后来加入了遗忘门来选择性地清除部分记忆。

\subsection{输出门的作用}

隐藏状态 $h^{(t)}$ 承担双重职责：（1）用于预测当前位置的下一个词；（2）存储可能对未来预测有用的上下文信息。输出门使得 LSTM 可以将这两个功能分离——细胞状态存储所有可能有用的信息，而输出门选择性地将当前需要的部分传递到隐藏状态用于输出预测。

例如，如果之前的文本提到了``the King of Prussia''，这个信息需要被记住以供未来使用，但在预测当前紧邻的下一个词时可能不相关。输出门可以让这部分信息留在细胞状态中而不暴露到隐藏状态。

\subsection{LSTM 也有助于缓解梯度爆炸}

LSTM 的加法更新不仅解决了梯度消失，也有助于缓解梯度爆炸。因为梯度不再需要经过连续的矩阵乘法链——加法操作提供了一条``梯度高速公路''，使梯度信号可以更稳定地流动。

\subsection{本章小结}

\begin{itemize}
    \item LSTM 通过引入\textbf{细胞状态}和\textbf{三个门}（遗忘门、输入门、输出门），实现了信息的选择性保留和更新
    \item \textbf{加法操作}是 LSTM 的核心设计——它允许梯度沿细胞状态直接流动，解决了梯度消失问题
    \item 当遗忘门为 1、输入门为 0 时，信息可以无损地沿时间线性传递
    \item LSTM 使有效上下文窗口从 $\sim$7 个 token 扩展到数十甚至数百个 token
\end{itemize}

%% ========== 其他 RNN 架构 ==========

